%% quantmsdiann manuscript — Molecular & Cellular Proteomics (MCP)
%% Article type: Software
%%
%% Build: cd paper && make pdf
%% Figures: MCP column widths via mcpfigures.sty; SVG sources in ../analysis/figures/
%%
%% TODO before submission — author-fillable placeholders:
%%   methods.tex:  UP000005640 UniProt release (e.g. 2024_05)
%%   methods.tex:  Cell-line panel deposit count (N3) and instrument-model count (I)
%%   methods.tex:  Zenodo DOI for consolidated panel SDRF
%%   methods.tex:  AWS r6i instance specs (confirm family / generation)
%%   results.tex:  Peak memory cap (GB) on consolidation step + max raw-file count tested
%%   results.tex:  CPU-hours ratio between 10-node and 300-node runs (from trace.txt)
%%   figures.tex (Fig. 4 caption): Zenodo DOI for consolidated pan-cohort SDRF + configurations
%%   declarations.tex: Competing-interest, Funding, Author contributions, Acknowledgements (currently draft placeholders)
%%
%% Pipeline version: v2.1.0 (release codename "Sao Pablo", 2026-05-08) — filled in methods.tex and declarations.tex

\documentclass[preprint,12pt]{elsarticle}
%% For journal submission, switch to double-column layout, e.g.:
%% \documentclass[final,3p,times,twocolumn]{elsarticle}

\usepackage[utf8]{inputenc}
\usepackage[T1]{fontenc}
\usepackage{microtype}
\usepackage{grffile}
\usepackage{amsmath}
\usepackage{url}
\usepackage{hyperref}
\usepackage{mcpfigures}
%% Shared macros for draft placeholders.
\newcommand{\placeholder}[1]{\textit{[#1]}}
\newcommand{\diannversion}[1]{\texttt{DIA-NN}~#1}
\newcommand{\quantmsdiann}{\texttt{quantmsdiann}}
\newcommand{\qpx}{\textsc{qpx}}


\journal{Molecular \& Cellular Proteomics}

\begin{document}

\begin{frontmatter}

\title{quantmsdiann: a scalable SDRF-driven DIA-NN workflow for reanalysis of single-cell, spatial, and bulk proteomics datasets}

\author[aff1]{Qi-Xuan Yue}
\author[aff1]{Yufei Shen}
\author[aff3]{Asier Larrea}
\author[aff2]{Chengxin Dai}
\author[aff4]{Henry Webel}
\author[aff6]{Jose Nimo}
\author[aff6]{Fabian Coscia}
\author[aff7]{Orhun Kok}
\author[aff7,aff8]{Nikolai Slavov}
\author[aff1]{Mingze Bai}
\author[aff5]{Vadim Demichev}
\author[aff3]{Yasset Perez-Riverol\corref{cor1}}

\ead[cor1]{yasset@ebi.ac.uk}
\cortext[cor1]{Corresponding author.}

\affiliation[aff1]{organization={Chongqing Key Laboratory of Big Data for Bio Intelligence, Chongqing University of Posts and Telecommunications}, city={Chongqing}, country={China}}
\affiliation[aff2]{organization={State Key Laboratory of Proteomics, Beijing Proteome Research Center, National Center for Protein Sciences (Beijing), Beijing Institute of Life Omics}, city={Beijing}, country={China}}
\affiliation[aff3]{organization={European Molecular Biology Laboratory, European Bioinformatics Institute (EMBL-EBI), Wellcome Genome Campus}, city={Hinxton, Cambridge}, country={UK}}
\affiliation[aff4]{organization={Novo Nordisk Foundation Center for Protein Research, Faculty of Health and Medical Sciences, University of Copenhagen}, city={Copenhagen}, country={Denmark}}
\affiliation[aff5]{organization={Quantitative Proteomics laboratory, Charit\'{e} -- Universit\"{a}tsmedizin Berlin}, city={Berlin}, country={Germany}}
\affiliation[aff6]{organization={Spatial Proteomics Group, Max-Delbr\"{u}ck-Center for Molecular Medicine in the Helmholtz Association (MDC)}, city={Berlin}, country={Germany}}
\affiliation[aff7]{organization={Department of Bioengineering, Department of Biology, Department of Chemistry and Chemical Biology, Single Cell Proteomics Center, and Barnett Institute for Chemical and Biological Analysis, Northeastern University}, city={Boston, Massachusetts}, country={USA}}
\affiliation[aff8]{organization={Parallel Squared Technology Institute}, city={Watertown, Massachusetts}, country={USA}}

\begin{abstract}
Public proteomics archives now host thousands of data-independent acquisition (DIA) datasets, but they remain analytically siloed because each was processed with a different software configuration. We present \quantmsdiann{}, an open-source Nextflow nf-core workflow that parallelize \texttt{DIA-NN} tool in cloud and HPC infrastructures guided by the experimental design written in SDRF format. The workflow ships container-pinned profiles for any \texttt{DIA-NN} version using Docker and Singularity packaging systems without the need of installations; dependency resolver etc. quantmsdiann support natively all vendor-formats and when conversion is needed (e.g. Sciex) the workflow handles it seamlessly the conversion to mzML. Results are exports harmonised quantification tables into MSstats input, a the newly Quantiative Proteomics eXchange format (\qpx{}), and a \texttt{pmultiqc} quality-control report. The \quantmsdiann{} parallelisation design allows a 2{,}300-run single-cell dataset to be reanalysed in 2.4~hours on an HPC cluster of 300 nodes.

We benchmark \quantmsdiann{} on the full ProteoBench DIA panel (five \texttt{DIA-NN} versions $\times$ four instrument modalities), reanalyse three independent bulk cell-line cohorts (PXD003539, PXD004701, PXD030304), and assemble a pan-cohort of five public DIA cell-line reanalyses under a single SDRF-driven invocation; SDRF-driven reprocessing with current \texttt{DIA-NN} versions matches or exceeds the originally deposited results. Single-cell and spatial proteomics demonstrations are part of ongoing work.
\quantmsdiann{} is available at \url{https://github.com/bigbio/quantmsdiann}.
\end{abstract}

\begin{keyword}
DIA-NN \sep single-cell proteomics \sep spatial proteomics \sep plexDIA \sep data-independent acquisition \sep Nextflow \sep nf-core \sep SDRF \sep reproducibility \sep QPX \sep pmultiqc
\end{keyword}

\end{frontmatter}

\section{Introduction}
\label{sec:introduction}
Data-independent acquisition (DIA) is now the dominant mass-spectrometry strategy for proteomics at scale, and \texttt{DIA-NN}~\cite{demichev2020} is among the most widely adopted analysis engines across timsTOF~\cite{brunner2022,meier2020}, Astral, and Orbitrap platforms.
Recent benchmarks confirm \texttt{DIA-NN}'s competitiveness for single-cell DIA workflows~\cite{wang2025}, and recent releases have expanded support for plexDIA, timsTOF parallel accumulation-serial fragmentation (PASEF), and vendor-native raw-file reading.
Despite this analytical maturity, most \texttt{DIA-NN} users run the tool as a desktop application, which limits its scalability for large datasets. Deployments on HPC and cloud infrastructure remain improvised, relying on shell scripts, bespoke directory conventions, and hand-curated sample sheets. As a result, the same raw data analysed by two laboratories, or with two \texttt{DIA-NN} releases, can yield non-overlapping quantification tables that are difficult to reconcile.

Beyond per-study analysis, the value of public proteomics archives such as PRIDE Archive~\cite{perezriverol2025pride} and other ProteomeXchange partners~\cite{deutsch2026proteomexchange} lies as much in systematic reanalysis as in original deposition. In 2024 we published bigbio/quantms~\cite{dai2024quantms}, a Nextflow pipeline that standardised multi-engine DIA and DDA analyses behind a Sample and Data Relationship Format (SDRF)~\cite{dai2021sdrf} input on the nf-core framework~\cite{ewels2020nfcore}. quantms was, however, built around \texttt{DIA-NN}~1.8.1 with a fixed parameter set, and its DDA/DIA breadth made extending DIA-specific functionality difficult. New \texttt{DIA-NN} parameters conflicted with DDA options, versions were hard-coded, sample-level acquisition parameters were not properly propagated to per-run \texttt{DIA-NN} invocations, and raw-file conversion was mandatory. Iterative reanalysis at archive scale requires a workflow that combines the reproducibility of nf-core, the sensitivity of current \texttt{DIA-NN} releases, the metadata discipline of SDRF, and a parallelisation strategy capable of processing thousands of raw files.

Here, we present \quantmsdiann{}, an independent SDRF-driven Nextflow pipeline dedicated to \texttt{DIA-NN} that extends DIA-specific features beyond what the original bigbio/quantms pipeline could accommodate. While \quantmsdiann{} shares its SDRF-first design philosophy and parts of the configuration-generation tooling with quantms, its workflow, container set, and module library are developed and released independently, with container-pinned profiles for every \texttt{DIA-NN} version selectable at runtime under Docker or Singularity. Downstream, it exports a harmonised \qpx{} (Quantitative Proteomics eXchange) Parquet archive alongside the standard \texttt{DIA-NN} reports and MSstats input (Methods). We demonstrate the pipeline along three axes: ProteoBench~\cite{devreese2025proteobench} benchmarks across four DIA modalities and successive \texttt{DIA-NN} releases; independent reanalyses of public DIA deposits spanning bulk cell lines (NCI-60, ProCan), a single-cell oocyte plexDIA dataset, a spatial Deep Visual Proteomics cohort, and a phosphoproteomics cohort; and a pan-cohort of five public DIA cell-line reanalyses integrated under a single SDRF-driven invocation.


\section{Experimental Procedures}
\label{sec:methods}
\subsection{Pipeline implementation}
\label{sec:impl}

\quantmsdiann{} is implemented in Nextflow DSL2 ($\geq$25.10.4) and follows nf-core community standards~\cite{ewels2020nfcore} for pipeline structure, testing, documentation, and continuous integration.
Each computational step is packaged as a versioned container image (Docker, Singularity, or Apptainer); upstream community tools (\texttt{DIA-NN}~1.8.1, ThermoRawFileParser~\cite{hulstaert2020}, pridepy~\cite{kamatchinathan2025pridepy}, qpx) are pulled from BioContainers, while pipeline-specific images (newer \texttt{DIA-NN} releases and \texttt{wiffconverter}) are hosted on \texttt{ghcr.io/bigbio}.
The repository at \url{https://github.com/bigbio/quantmsdiann} is organised as a standard nf-core pipeline with main workflow (\texttt{main.nf}), subworkflows, modules, test data and CI configurations, and documentation.
Pipeline version v2.1.0 (release codename ``Sao Pablo'', 2026-05-08) was used for the analyses reported in this paper.

\subsection{SDRF parsing and configuration generation}
\quantmsdiann{} accepts SDRF~\cite{dai2021sdrf} as primary input.
SDRF parsing is performed using the sdrf-pipelines Python package developed within the quantms project, which validates the SDRF against the proteomics-specific subset of EFO, PSI-MS, and PRIDE ontology terms and emits a normalised per-sample channel for downstream Nextflow processes.
\texttt{DIA-NN} configuration generation extracts enzyme, fixed and variable modification, instrument, precursor mass tolerance, fragment mass tolerance, and missed-cleavage settings from the SDRF and writes a \texttt{DIA-NN} command-line configuration.
Where SDRF columns are absent, \texttt{DIA-NN} defaults appropriate for the inferred instrument are applied.

\subsection{Raw file handling and \texttt{DIA-NN} execution}
Raw inputs are dispatched to format-specific handlers in the file-preparation subworkflow, each conversion or decompression running as an independent, parallelised Nextflow process.
Thermo \texttt{.raw} files are read natively by \texttt{DIA-NN}~$\geq$~2.1.0 (which added Linux native-reader support) or converted to mzML with ThermoRawFileParser~\cite{hulstaert2020} otherwise (\texttt{--mzml\_convert}); Bruker \texttt{.d} folders are ingested natively, with compressed forms (\texttt{.d.tar}/\texttt{.zip}/\texttt{.gz}) decompressed first; and Sciex \texttt{.wiff}~/~\texttt{.wiff2} are converted to mzML with the open-source \texttt{wiffconverter} (\url{https://github.com/bigbio/wiffconverter}), removing the historical dependency on a proprietary ProteoWizard installation.
Generic \texttt{.dia} and pre-converted \texttt{.mzML} inputs pass through unchanged, with optional OpenMS re-indexing (\texttt{--reindex\_mzml}).
\texttt{DIA-NN}~\cite{demichev2020} is then invoked at the version chosen via \texttt{--diann\_version}: the ProteoBench analyses use \diannversion{1.8.1}, \diannversion{2.1.0}, \diannversion{2.2.0}, \diannversion{2.3.2}, and \diannversion{2.5.0} from the same SDRF, and the cell-line reanalyses use \diannversion{2.5.0} unless otherwise stated.

\subsection{Spectral-library generation}
\quantmsdiann{} runs \texttt{DIA-NN} library-free, generating the spectral library in silico from the input FASTA rather than supplying an external experimental library: \texttt{DIA-NN}'s deep-learning predictor builds a predicted library from the protein sequences, each run is searched against it, a cohort-wide empirical library is assembled from the confident first-pass identifications, and a final pass re-searches and quantifies against that refined library.
The source FASTA was the UniProt human reference proteome~\cite{uniprot2025} (UP000005640, \placeholder{release}), with match-between-runs enabled at the cohort level.
In ProteoBench's submission taxonomy this is the predicted-library strategy (\texttt{predictors\_library}~$=$~DIANN); the ProteoBench comparison (Section~\ref{sec:proteobench}) is therefore restricted to community submissions that used the same predicted-from-FASTA strategy.

\subsection{Quality control and \qpx{} export}
Per-run and per-cohort quality control reports are generated by \texttt{pmultiqc}~\cite{larrea2025pmultiqc}, a proteomics extension of MultiQC~\cite{ewels2016multiqc} that summarises \texttt{DIA-NN}-specific metrics including peptide and protein counts, missed cleavages, charge-state distributions, precursor and fragment mass accuracy, and per-run identification yield.
The \qpx{} archive is produced by the \texttt{qpxc} CLI from the \texttt{qpx} package (\url{https://github.com/bigbio/qpx}), which writes Parquet files for identifications (psm, peptide, protein-group, feature), metadata (sample, run, ontology), and provenance, and optionally exports MuData (\texttt{.h5mu}) containers for scverse-based~\cite{virshup2023scverse} downstream analysis.

\subsection{Datasets benchmarked}
\label{sec:datasets}
The benchmark and reanalysis datasets used in this study (ProteoBench DIA modules~5/7/9/10, the three per-cohort cell-line reanalyses PXD003539, PXD004701, PXD030304, the two additional pan-cohort datasets PXD041421 and PXD017199, and the PXD071075 scalability sweep) are summarised together with their instrument, acquisition mode, and deposition year in Additional file~1: Table~S1.
For each dataset, an SDRF was prepared from the original sample metadata and is provided as part of the Additional file~1 deposition.
For each cell-line reanalysis cohort, the original published quantification was downloaded from the PRIDE deposition and used as the comparator, thresholded at 1\% precursor-level FDR for both \quantmsdiann{} and the original analyses.
The conservative contaminant/decoy filter and the like-for-like-threshold rule applied to the original-paper headline counts are described in Additional file~1.

\subsection{Compute environment and benchmarking}
Runtime benchmarks were performed on the EMBL-EBI HPC cluster (LSF scheduler), with per-task CPU and memory allocations dispatched by Nextflow.
Memory profiling was performed using the Nextflow trace report (\texttt{-with-trace}).



\section{Results}
\label{sec:results}
\subsection{\quantmsdiann{} pipeline architecture}
\label{sec:architecture}

Built on the nf-core framework~\cite{ewels2020nfcore}, \quantmsdiann{} accepts a single SDRF file as primary input, which encodes per-sample metadata (cell line or tissue origin, factor values, instrument, post-translational modifications) and per-run file references. We extended the SDRF with DIA-relevant columns (MS1 and MS2 mass windows, and plexDIA channel definitions and labels) and added an adapter to sdrf-pipelines~\cite{dai2021sdrf} (\texttt{convert-diann}) that translates each run's declared metadata into a per-run \texttt{DIA-NN} command-line configuration. Every run is therefore parametrised from its own metadata with no manual configuration or parameter-harmonisation step, removing a common source of analysis-script divergence between groups: different MS1/MS2 window schemes can be mixed within a cohort, or a plexDIA dataset run with its channel definitions, without any workflow change beyond the SDRF.
The workflow comprises four mandatory and three optional modules organised around the \texttt{DIA-NN} analysis core (Fig.~\ref{fig:architecture}a), alternating between per-file parallel stages (red; raw-file preparation, preliminary calibration, and individual search) and collective stages (green; in-silico library generation, empirical library assembly, final quantification, MSstats conversion, and \texttt{pmultiqc}~\cite{larrea2025pmultiqc} quality-control reporting) operating on the full cohort.

\quantmsdiann{} follows \texttt{DIA-NN}'s distributed multi-pass flow, alternating parallel per-file work with serial cohort-wide steps (Fig.~\ref{fig:architecture}a). Each raw file is first prepared \emph{in parallel} (vendor-format conversion, decompression, or indexing) where needed; conversion to mzML is mandatory for early \texttt{DIA-NN} versions and for some vendor formats such as SCIEX \texttt{.wiff}. A spectral library is then predicted in-silico from the FASTA once (\emph{serial}; skipped when a library is supplied), and every file is calibrated against it in a \emph{parallel} first pass. The first-pass results are merged into a single empirical library (\emph{serial}); each file is searched again against that library \emph{in parallel}; and a final consolidation pass merges all per-file results into the cohort quantification, applying match-between-runs, normalisation, and protein inference (\emph{serial}). MSstats conversion, \texttt{pmultiqc} reporting, and \qpx{} export then run once at the end (\emph{serial}). The two per-file search passes carry essentially all of the compute and scale across cluster nodes, driving the throughput reported below (Section~\ref{sec:throughput}); the serial steps set a largely fixed wall-clock floor.
Outputs include the standard \texttt{DIA-NN} report tables: \texttt{report.tsv} (or \texttt{report.parquet}), \texttt{report.pr\_matrix.tsv}, and \texttt{report.pg\_matrix.tsv}, alongside MSstats-formatted input and a \texttt{pmultiqc}~\cite{larrea2025pmultiqc} quality-control report.
A harmonised \qpx{} archive bundles identifications, metadata, and provenance into a Parquet-based artefact, queried directly with DuckDB and, for the quantification layers, exported to MuData (\texttt{.h5mu}) for scverse-compatible~\cite{virshup2023scverse} downstream analysis.

\subsection{quantmsdiann throughput and scaling}
\label{sec:throughput}

To characterise how \quantmsdiann{} scales on an HPC cluster, we simulated the reanalysis of an entire single-cell experiment of \textasciitilde 2{,}300 raw files (PXD071075) across cluster sizes from 10 to 300 nodes (set through the Nextflow \texttt{executor.queueSize}), and measured the end-to-end wall-clock time (Fig.~\ref{fig:architecture}b). Because the per-file work is parallelised across nodes, a single SDRF-driven invocation completed the whole cohort within hours: wall-clock fell monotonically from 37.4~h at 10 nodes to 2.2~h at 300 nodes (excluding the one-time in-silico library-prediction step, which is independent of cluster size). The curve has a practical knee near 200 nodes (2.4~h), beyond which the serial cohort-level steps dominate and additional nodes give diminishing returns (a 30$\times$ increase in cluster width yields only a 17.4$\times$ speed-up), so the cluster should be sized to the desired turn-around time rather than to maximum throughput.

We then compared the time-to-finish across every dataset in this study (Fig.~\ref{fig:architecture}c), reporting each run net of the one-time in-silico library-prediction step (a cohort-independent cost that can be predicted once and reused). For proteome-scale cohorts, wall-clock stayed within a few hours irrespective of cohort size, because the per-file passes absorb additional files in parallel: the largest cohort (ProCan, PXD030304; 5{,}798 TripleTOF~6600 runs) finished in 6.2~h, in the same few-hour band as cohorts orders of magnitude smaller, while the 6-run ProteoBench benchmarks completed in under 1.5~h. Time-to-finish is therefore governed by the per-file search depth and the remaining cohort-level consolidation rather than by the number of files: the single longest run is a 10-file phosphoproteomics cohort (PXD049692, 8.0~h after excluding library prediction), where the deep variable-modification search dominates the residual wall-clock. Peak memory during the cohort-level consolidation stayed within the 64--128~GB available on commodity HPC nodes, so no high-memory specialist hardware is required. Per-step wall-clock and resource envelopes are reported in Additional file~1: Figs.~S1 and~S2.

\subsection{ProteoBench benchmark across \texttt{DIA-NN} versions and instruments}
\label{sec:proteobench}

To test whether \quantmsdiann{} preserves \texttt{DIA-NN}'s analytical performance under SDRF-driven orchestration, we processed every public ProteoBench DIA module through the workflow and compared the resulting identification counts against the ProteoBench public-submission snapshot (29 May 2026)~\cite{devreese2025proteobench}.
Because \quantmsdiann{} runs \texttt{DIA-NN} library-free (predicting the library in-silico from the FASTA, with no externally supplied experimental library), we restricted the community set for a like-for-like comparison to ProteoBench \texttt{DIA-NN} submissions that used the same predicted-from-FASTA library (\texttt{predictors\_library}~$=$~DIANN); submissions that loaded a pre-existing experimental or user-curated spectral library search a different candidate space and are excluded from the headline comparison.
The benchmark panel covers four instrument families and acquisition modes: Orbitrap Astral 2~Th DIA (\emph{Module~7}), single-cell DIA on Orbitrap Astral (\emph{Module~9}, PXD049412), timsTOF diaPASEF (\emph{Module~5}, PXD062685), and SCIEX ZenoTOF~7600 (\emph{Module~10}, PXD070049).
For each module we ran the workflow against three \texttt{DIA-NN} configurations --- the oldest supported release (1.8.1) and the current release in its academic (2.5.1) and enterprise (2.5.1-enterprise) builds --- in a single invocation, all parametrised from the same SDRF (Fig.~\ref{fig:proteobench}). Identification counts are taken from the per-precursor \texttt{DIA-NN} report rather than the quantities matrices. Each release is run at \texttt{DIA-NN}'s recommended run-specific precursor q-value (\texttt{--qvalue}: 1\% for 1.8.1, 5\% for the 2.5.1 builds); on top of this we apply, uniformly across configurations, a 1\% global precursor q-value, and count unique precursors identified in at least one run. Protein groups are summarised --- like-for-like across versions, each at a 1\% run-specific protein-group q-value --- as the per-run average and the count quantified in \emph{every} run (``complete profiles''); the per-run depth metric tracks a release's gains more directly than the cross-run union, which saturates as rare identifications accumulate across replicates. The recommended run-specific precursor cut-off mainly affects the reproducibility-filtered ($\geq$3-replicate) precursor counts; the single-replicate precursor totals and the per-run protein metrics apply the same thresholds to every version and are directly comparable across releases.

Precursor counts (identified in at least one run) rose from 1.8.1 to 2.5.1 on every module, by 12\% on Module~7 (116{,}771~$\to$~130{,}246), 11\% on Module~9 (41{,}756~$\to$~46{,}243), 21\% on Module~5 (97{,}149~$\to$~117{,}396), and 16\% on Module~10 (88{,}179~$\to$~102{,}457); the enterprise build adds a further 2--3\% on each module.
Protein-group depth rose in step. Averaged per run, protein groups at 1\% FDR increased from 1.8.1 to 2.5.1 by 8--11\% (Module~7 10{,}091~$\to$~11{,}167; Module~9 7{,}157~$\to$~7{,}895; Module~5 10{,}199~$\to$~11{,}231; Module~10 8{,}788~$\to$~9{,}470), with the enterprise build a further 1--3\% higher and highest on every module. The improvement is larger for proteins quantified in \emph{every} run (complete profiles): the enterprise build gains 12--20\% over 1.8.1 (Module~7 8{,}799~$\to$~10{,}403, $+$18\%; Module~9 6{,}124~$\to$~7{,}349, $+$20\%; complete-profile panel in Additional file~1).
Reproducibility-filtered precursor counts ($\geq$3 replicates) follow the same ordering across the three configurations on all four modules.
This monotonic version-progress trend is consistent with what the broader ProteoBench community reports across \texttt{DIA-NN} releases~\cite{devreese2025proteobench}, indicating that the SDRF wrapper introduces no version-coupling artefacts.
Quantification accuracy, by contrast, is essentially invariant across \texttt{DIA-NN} versions.
On every module the measured per-species log$_2$ fold-changes recover the expected HYE ratios (human at unity, yeast and \emph{E.\ coli} near their design ratios, with the mild attenuation at the extreme \emph{E.\ coli} ratio that is characteristic of DIA quantification), and the three \texttt{DIA-NN} configurations overlie one another at every species (Fig.~\ref{fig:proteobench}b).
The version-to-version shift in median quantification error~$|\epsilon|$ ($\sim$0.01 log$_2$, $\approx$0.5--1\% in linear fold-change) is far smaller than the per-precursor spread, so the choice of \texttt{DIA-NN} release is best driven by the identification gains in panel~(a) rather than by accuracy.
On the two modules for which the snapshot contained predicted-library \texttt{DIA-NN} submissions, Orbitrap Astral (Module~7, $n=15$) and timsTOF diaPASEF (Module~5, $n=8$), \quantmsdiann{}'s three configurations fall within the community median-$|\epsilon|$ distribution, at or below its median (Fig.~\ref{fig:proteobench}c); the single-cell and ZenoTOF modules carried no predicted-library community submission. Per-version, per-module protein-group and precursor breakdowns are provided in Additional file~1: Fig.~S3.

\subsection{Independent reanalyses of public DIA deposits recover comparable or deeper coverage than the original quantifications}
\label{sec:reanalyses}

To evaluate \quantmsdiann{} on production-scale public deposits, we reanalysed independent DIA cohorts, spanning bulk cell lines, single cells, and spatial proteomics, that had been processed and published with different DIA engines (Fig.~\ref{fig:cell-line-reanalysis}).
For each cohort the same raw files were retrieved from PRIDE Archive~\cite{deutsch2026proteomexchange} (or MassIVE, for the single-cell plexDIA cohort), parsed through an SDRF generated from the original sample annotations, and run through \quantmsdiann{} with current \texttt{DIA-NN} releases.

On the \textbf{NCI-60} panel (PXD003539, Guo \textit{et al.}~2019~\cite{guo2019nci60}; originally analysed with OpenSWATH against the Guo assay library), \quantmsdiann{} recovered 95{,}633 peptides at 1\% FDR against 40{,}592 in the original submission ($+135$\%), with comparable protein-group counts (6{,}747 versus 6{,}556; Fig.~\ref{fig:cell-line-reanalysis}a).
On the \textbf{single-cell oocyte plexDIA} cohort (MSV000093870, Galatidou \textit{et al.}~2024~\cite{galatidou2024oocyte}; mTRAQ 3-channel, Q~Exactive), the first plexDIA dataset processed through \quantmsdiann{}, the workflow quantified 2{,}904 protein groups from the same raw data, versus 2{,}122 in the original analysis ($+37$\%). At the protein-group level it recovered 94\% of the originally reported groups (1{,}999 of 2{,}122) and detected 905 further groups, at comparable per-cell median depth (1{,}736 versus 1{,}784, across 107 and 97 cells respectively; Fig.~\ref{fig:cell-line-reanalysis}b). A QC-matched per-oocyte comparison (Additional file~1) confirms equivalent single-cell depth (Pearson $r=0.98$).
On the \textbf{ProCan}-DepMapSanger 949-cell-line pan-cancer panel (PXD030304, Gon\c{c}alves \textit{et al.}~2022~\cite{goncalves2022procan}; originally analysed with \texttt{DIA-NN}~1.7.x), \quantmsdiann{} recovered 8{,}921 protein groups at 1\% FDR against 8{,}498 originally ($+5$\%) and 7{,}990 protein groups requiring $\geq 2$ unique peptides versus 6{,}692 ($+19$\%; Fig.~\ref{fig:cell-line-reanalysis}c).
Both original counts are reported in the ProCan paper itself, and the \quantmsdiann{} numbers use the matching ProCan-paper filters under our conservative contaminant/decoy policy (Experimental Procedures), so each side of the comparison is like-for-like.
A fourth, bulk cohort already analysed with a current DIA engine, the Sun \textit{et al.}~2023 PCT-SWATH breast-cancer panel (PXD004701, 76 lines~\cite{sun2023breast}), yielded near-parity (6{,}146 versus 6{,}091 protein groups, $+0.9$\%; Additional file~1), as expected when the original analysis already used a modern engine.
On a \textbf{spatial} Deep Visual Proteomics cohort of MYCN-amplified neuroblastoma (PXD064049; laser-microdissected regions, diaPASEF, originally analysed with \texttt{DIA-NN}~1.8.1), \quantmsdiann{} (\texttt{DIA-NN}~2.5.0) matched the original at the precursor level (16{,}443 versus 16{,}518) and quantified somewhat fewer protein groups (2{,}425 versus 2{,}882; Fig.~\ref{fig:cell-line-reanalysis}d); this protein-group difference reflects \quantmsdiann{}'s entrapment-augmented search, which imposes tighter error control than the original plain-FASTA analysis (only 0.7\% of \quantmsdiann{}'s accepted protein groups map to entrapment sequences), rather than reduced sensitivity.
Finally, on a \textbf{phosphoproteomics} cohort (PXD049692; NK-cell Fe-NTA enrichment, diaPASEF, originally analysed with Spectronaut directDIA), \quantmsdiann{} recovered a comparable number of phosphopeptide backbones at 1\% FDR (distinct stripped sequences: 4{,}234 versus 4{,}254; Fig.~\ref{fig:cell-line-reanalysis}e). At the site-resolved phosphopeptidoform level the two analyses diverge (4{,}592 versus 7{,}993), with Spectronaut reporting more positional isomers; we therefore frame this as comparable phosphopeptide-backbone coverage rather than site-level parity.
Per-condition completeness, peptide-per-protein distributions, gene-versus-accession overlap and per-cell-line missingness for the cell-line and single-cell cohorts are provided as Additional file~1: Figs.~S4--S12.

\subsection{Pan-cohort of five public DIA cell-line reanalyses}
\label{sec:pancohort}

To illustrate the value of running the same SDRF-driven workflow across many independent deposits, we combined the three reanalysed cohorts above with two additional cell-line DIA panels (PXD041421, the MultiPro batch-effect testbed, Wang \textit{et al.}~2023~\cite{wang2023multipro}; PXD017199, Tognetti \textit{et al.}~2021~\cite{tognetti2021breast}, 67 breast-cancer lines) into a single \quantmsdiann{} invocation and assembled a pan-cohort view of per-cohort yield and per-tissue coverage (Fig.~\ref{fig:pancohort}).
Per-cohort headline counts are summarised in Fig.~\ref{fig:pancohort}a.
Across the five cohorts the reanalysis spans 12{,}229 unique UniProt accessions, of which 5{,}144 (42.1\%) are detected in all five (a pan-cohort core proteome) and a further 1{,}334 (10.9\%) in exactly four; these overlap counts are tabulated in the combined per-cohort counts provided with the analysis code.

Tissue-level coverage across the combined panel spans 27 Cellosaurus-derived tissue categories (Fig.~\ref{fig:pancohort}b,c): lung, breast, haematopoietic-and-lymphoid, skin, and central nervous system together account for 59\% of all cell-line entries in the pan-cohort.
Because every count comes from a single SDRF-parameterised \quantmsdiann{} invocation, the resulting protein matrix is queryable from the published \qpx{} archive and can be filtered, joined, or re-grouped by Cellosaurus, NCIt, or Disease Ontology terms without manual harmonisation.


%% Figure environments for the quantmsdiann MCP manuscript.

%% Fig. 1 — pipeline architecture (subway, vector PDF from rsvg-convert)
%% on the left, runtime scaling panels stacked on the right.
\begin{figure}[htbp]
  \centering
  \begin{minipage}{\textwidth}
    \centering
    \begin{minipage}[t]{0.62\linewidth}
      \centering
      \textbf{(a)}\par\vspace{0.2em}
      \mcpincludegraphics[width=\linewidth,keepaspectratio]{manuscript/workflow_panel}
    \end{minipage}\hfill
    \begin{minipage}[t]{0.36\linewidth}
      \centering
      \textbf{(b)}\par\vspace{0.2em}
      \mcpincludegraphics[width=\linewidth,keepaspectratio]{performance/queue_size_sweep}\par\vspace{1.0em}
      \textbf{(c)}\par\vspace{0.2em}
      \mcpincludegraphics[width=\linewidth,keepaspectratio]{performance/parallelism_vs_wallclock}
    \end{minipage}
  \end{minipage}
  \caption[Pipeline and scaling]{%
    \textbf{\quantmsdiann{} pipeline architecture and runtime scaling.}
    (\textbf{a})~Subway diagram of the \quantmsdiann{} workflow: mandatory
    modules as filled circles, optional/skippable modules as open circles.
    Per-file parallel stages (red; file preparation, preliminary
    calibration, individual search) run as one task per raw file, whereas
    collective stages (green; in-silico library generation, empirical
    library assembly, final quantification, MSstats conversion, and
    \texttt{pmultiqc} reporting) run once over the full cohort.
    (\textbf{b})~Workflow wall-clock time versus the number of cluster nodes
    (Nextflow \texttt{executor.queueSize}) on the PXD071075 single-cell
    cohort.
    (\textbf{c})~Wall-clock time to finish each reanalysis, one bar per
    dataset, ordered by cohort size (number of MS data files) and coloured
    by instrument family. In both~(\textbf{b}) and~(\textbf{c}) the reported
    time \emph{excludes} the in-silico spectral-library prediction step
    (\texttt{INSILICO\_LIBRARY\_GENERATION}), a one-time, cohort-independent
    cost that can be predicted once and reused, so the bars reflect the
    per-file quantification work. Time-to-finish stays within a few hours
    across two orders of magnitude in file count because per-file stages are
    parallelised across the cluster; the deeper variable-modification search
    of the phosphoproteomics cohorts (e.g.\ PXD049692) is the main residual
    driver of wall-clock once library prediction is set aside.
  }
  \label{fig:architecture}
\end{figure}

%% Fig. 2 — ProteoBench: precursors-vs-version (compact summary) +
%% ID-vs-accuracy scatter (community context). The per-version per-module
%% protein-group / precursor breakdown is in Additional file 1 (Fig. S3).
\begin{figure}[htbp]
  \centering
  \mcpincludegraphics[width=\linewidth,keepaspectratio]{manuscript/fig2_proteobench_equivalence}
  \caption[quantmsdiann reproduces single-machine ProteoBench submissions]{%
    \textbf{Distributed \quantmsdiann{} reproduces single-machine results on
    ProteoBench.} For the two DIA modules with public predicted-from-FASTA
    community submissions (Orbitrap Astral, Module~7; timsTOF diaPASEF,
    PXD062685), \quantmsdiann{} (\texttt{DIA-NN} 1.8.1 and 2.5.1-enterprise) is
    compared against the cloud of community submissions, which are
    single-machine sequential runs. \quantmsdiann{} preserves
    match-between-runs by design: a preliminary pass builds a combined
    empirical library across all runs, against which each run is re-quantified.
    (\textbf{a})~Accuracy per concentration --- measured log$_2$ fold-change for
    each HYE species against the ProteoBench-expected ratio (dashed); community
    submissions as grey points.
    (\textbf{b})~Depth versus accuracy --- precursors quantified against
    median-$|\epsilon|$; \quantmsdiann{} (coloured) sits within the
    single-machine community cloud (grey), reaching greater depth at comparable
    accuracy from 1.8.1 to 2.5.1-enterprise.
    Datasets without predicted-library community comparators (single-cell
    Module~9/PXD049412; ZenoTOF Module~10/PXD070049) and the per-version
    protein-group counts are in Additional file~1.
  }
  \label{fig:proteobench}
\end{figure}

%% Fig. 3 — Independent cell-line reanalyses (depth vs original)
\begin{figure}[htbp]
  \centering
  \begin{minipage}{\textwidth}
    \centering
    \begin{minipage}[t]{0.32\linewidth}
      \centering
      {\scriptsize\textbf{(a)~PXD003539 (NCI-60, 2019)}}\par\vspace{0.2em}
      \mcpincludegraphics[width=\linewidth,keepaspectratio]{PXD003539/main_comparison}
    \end{minipage}\hfill
    \begin{minipage}[t]{0.32\linewidth}
      \centering
      {\scriptsize\textbf{(b)~MSV000093870 (plexDIA, 2024)}}\par\vspace{0.2em}
      \mcpincludegraphics[width=\linewidth,keepaspectratio]{plexDIA/MSV000093870/main_galatidou_total_pg}
    \end{minipage}\hfill
    \begin{minipage}[t]{0.32\linewidth}
      \centering
      {\scriptsize\textbf{(c)~PXD030304 (ProCan, 2022)}}\par\vspace{0.2em}
      \mcpincludegraphics[width=\linewidth,keepaspectratio]{PXD030304/main_comparison}
    \end{minipage}

    \vspace{1.0em}
    \begin{minipage}[t]{0.32\linewidth}
      \centering
      {\scriptsize\textbf{(d)~PXD064049 (MYCN DVP spatial, 2025)}}\par\vspace{0.2em}
      \mcpincludegraphics[width=\linewidth,keepaspectratio]{PXD064049/main_comparison}
    \end{minipage}\hfill
    \begin{minipage}[t]{0.32\linewidth}
      \centering
      {\scriptsize\textbf{(e)~PXD049692 (NK phospho, 2024)}}\par\vspace{0.2em}
      \mcpincludegraphics[width=\linewidth,keepaspectratio]{PXD049692/main_comparison}
    \end{minipage}\hfill
    \begin{minipage}[t]{0.32\linewidth}~\end{minipage}
  \end{minipage}
  \caption[Cell-line, spatial, and phospho reanalyses]{%
    \textbf{Independent reanalysis of public DIA datasets, spanning bulk cell
    lines, single cells, spatial proteomics, and phosphoproteomics, recovers
    comparable or greater identification depth than the original deposits.}
    Each panel compares the original published counts (grey) with the
    \quantmsdiann{} reanalysis at 1\% FDR (blue) from the same raw data.
    (\textbf{a})~NCI-60 bulk cell lines (PXD003539, Guo 2019; original
    OpenSWATH analysis).
    (\textbf{b})~Single-cell oocyte plexDIA (MSV000093870, Galatidou 2024).
    (\textbf{c})~ProCan-DepMapSanger pan-cancer cell-line panel (PXD030304).
    (\textbf{d})~MYCN-amplified neuroblastoma Deep Visual Proteomics
    (spatial, laser-microdissected regions; PXD064049, diaPASEF), with
    precursors and protein groups compared against the original
    DIA-NN~1.8.1 analysis.
    (\textbf{e})~NK-cell Fe-NTA phosphoproteomics (PXD049692, diaPASEF):
    stripped-sequence phosphopeptides (backbones) compared against the
    original Spectronaut directDIA analysis on the same runs.
    Per-cohort detail for panels~(a)--(c) in Additional file~1: Figs.~S4--S12.
  }
  \label{fig:cell-line-reanalysis}
\end{figure}

%% Fig. 4 — Pan-cohort of five DIA reanalyses
\begin{figure}[htbp]
  \centering
  \mcpincludegraphics[width=\textwidth,height=\mcpmaxfigheight,keepaspectratio]{manuscript/fig4_cellline_atlas}
  \caption[Pan-cohort of DIA reanalyses]{%
    \textbf{Pan-cohort of five public DIA cell-line reanalyses by
    \quantmsdiann{}.}
    Combines PXD003539 (NCI-60, Guo 2019), PXD030304
    (ProCan-DepMapSanger, 949 lines), PXD004701 (Sun 2023, 76 breast
    cancer lines), PXD041421 (Wang 2023), and PXD017199 (Tognetti
    2021, 67 breast cancer lines), reanalysed under a single
    \quantmsdiann{} invocation.
    (\textbf{a})~Per-cohort headline protein-group counts: the originally
    published count (grey) and the \quantmsdiann{} reanalysis (blue);
    PXD017199 and PXD041421 have no published DIA headline, so only the
    \quantmsdiann{} bar is shown.
    (\textbf{b})~Cell lines per pan-cancer tissue, stacked by contributing
    cohort.
    (\textbf{c})~Unique UniProt proteins detected per tissue across the
    combined panel.
  }
  \label{fig:pancohort}
\end{figure}


\section{Discussion}
\label{sec:discussion}
\quantmsdiann{} lowers the engineering barrier between a deposited DIA dataset and a reproducible quantification table, and it does so at a scale that makes systematic reanalysis of public DIA archives tractable for the first time.
The SDRF-first design encourages experiment annotation upstream of analysis and produces metadata that is interoperable with the wider proteomics ecosystem~\cite{dai2021sdrf}; the \qpx{} archive bundles the SDRF, harmonised quantification tables and pipeline provenance into a single Parquet container so that reanalysed results can be archived (e.g.\ on Zenodo) and re-queried alongside the original raw data.
By focusing exclusively on \texttt{DIA-NN}, rather than the multi-engine scope of the related bigbio/quantms pipeline~\cite{dai2024quantms}, \quantmsdiann{} achieves a leaner dependency surface and faster pipeline-level execution for the DIA-only regime, where the dominant cost is per-file \texttt{DIA-NN} invocation rather than across-engine comparison.

The demonstrations we present stress the workflow along three orthogonal axes.
The ProteoBench benchmark (Fig.~\ref{fig:proteobench}) tests analytical fidelity across four instrument modalities and five \texttt{DIA-NN} releases; the three independent cell-line reanalyses (Fig.~\ref{fig:cell-line-reanalysis}) test recovery against published per-study quantifications, with results ranging from near-parity (Sun 2023, current DIA engine) to substantial recovery gains (Guo 2019 OpenSWATH; ProCan \texttt{DIA-NN}~1.7.x); and the pan-cohort of five public DIA cell-line reanalyses (Fig.~\ref{fig:pancohort}) tests archive-scale throughput on five concurrently reanalysed deposits.
Single-cell and spatial proteomics biology demonstrations are part of ongoing work and are not part of this submission's biology figures; the single-cell PXD071075 cohort used as the throughput-scaling benchmark in Fig.~\ref{fig:architecture} is, however, real production data.
Across the cell-line panels, the largest per-run gains accrue for deposits originally analysed with older or non-\texttt{DIA-NN} engines (Guo 2019 OpenSWATH, ProCan 1.7.x), and gains shrink to near-parity for the most recently processed cohorts (Sun 2023).
This argues that the value of DIA reanalysis is real but time-limited: it is largest for old data and shrinks for new data, which is exactly the regime an SDRF-driven workflow is well-positioned to address as new \texttt{DIA-NN} versions ship.
Demonstrations on flagship single-cell DIA and laser-microdissected spatial cohorts are part of the development roadmap; the workflow already supports these acquisition modes (Fig.~\ref{fig:architecture}a, library-handling branch) and SDRFs for the Brunner~\cite{brunner2022}, Derks~\cite{derks2023}, Leduc~\cite{leduc2024}, Wang~\cite{wang2025}, Mund~\cite{mund2022} and Makhmut~\cite{makhmut2023} datasets are in preparation.

The cross-\texttt{DIA-NN}-version reproducibility view enabled by \quantmsdiann{} is, to our knowledge, the first pipeline-level treatment of this question for \texttt{DIA-NN}.
We expect this to become increasingly important as \texttt{DIA-NN} evolves: a deposited single-cell DIA dataset analysed with \diannversion{1.8} today and re-analysed with \diannversion{2.x} in two years should not produce different biological conclusions for the same data and parameters, but in current practice such re-analyses are rarely performed because the bespoke per-study scripts make them prohibitively expensive.
\quantmsdiann{}'s container-pinned, version-parametrised execution makes this kind of audit a single-command operation.

Several limitations are worth noting.
First, \quantmsdiann{} does not currently include batch correction, cell-type assignment, or differential-abundance modules; these are best handled in downstream R or Python packages such as scp~\cite{vanderaa2023}, which can read the \qpx{} or MSstats outputs directly.
Second, library generation for plexDIA channel deconvolution remains internal to \texttt{DIA-NN}, and \quantmsdiann{} inherits any limitations of that step.
Third, the \qpx{} format will continue to evolve as community feedback accumulates; backward compatibility is maintained through schema versioning.
Fourth, we focused on bulk DIA cell-line reanalyses; plexDIA, TMT-DIA, and single-cell DIA demonstrations are part of the development roadmap but are not evaluated here.

Finally, we note that \quantmsdiann{} is complementary to, rather than competing with, existing single-cell DIA workflows.
The benchmarking work of Wang \textit{et al.}~\cite{wang2025} compared multiple search engines and rescoring strategies; \quantmsdiann{} adopts the broadly competitive \texttt{DIA-NN} engine and focuses on the orthogonal problem of orchestrating that engine reproducibly at scale.
Similarly, the recent MHCquant2 nf-core pipeline~\cite{scheid2025mhcquant2} addresses an analogous orchestration problem for immunopeptidomics.
Together, these contributions extend the nf-core proteomics ecosystem to cover the major MS-based proteomics modalities behind a common SDRF-driven interface.


\section*{Declarations}

\subsection*{Ethics approval and consent to participate}
Not applicable.
All datasets reanalysed in this study were obtained from public proteomics repositories with ethics approval and informed consent from the original depositing studies.

\subsection*{Consent for publication}
Not applicable.

\subsection*{Data availability}
All datasets reanalysed in this study are publicly available through ProteoMeXchange.
Cell-line bulk DIA reanalyses: PRIDE PXD003539 (Guo 2019, NCI-60), PXD004701 (Sun 2023), PXD030304 (ProCan-DepMapSanger), PXD017199 (Tognetti 2021), PXD041421 (Wang 2023).
Scalability experiment: PRIDE PXD071075.
ProteoBench benchmarks: PXD049412 (single-cell, Module 9), PXD062685 (diaPASEF, Module 5), PXD070049 (ZenoTOF, Module 10), and Module 7 (Astral) datasets.
The full list of PRIDE accessions used in the pan-cohort of DIA cell-line reanalyses is provided in Additional file~1: Table~S1.
Per-process runtime and resource envelopes (Figs.~S1, S2) and detailed ProteoBench supplementary panels (Figs.~S3, S4) are provided as Additional file~1.

\subsection*{Code availability}
\quantmsdiann{} is available at \url{https://github.com/bigbio/quantmsdiann} under the MIT license; documentation is hosted at \url{https://quantmsdiann.quantms.org/}.
The pipeline is implemented in Nextflow DSL2 ($\geq$25.10.4) and packages each step as a versioned Docker, Singularity, or Apptainer container image.
\texttt{DIA-NN}~1.8.1 is distributed publicly via BioContainers; because the \texttt{DIA-NN} license restricts redistribution of later releases, containers for \texttt{DIA-NN}~1.9 and above are built locally from the recipes in \url{https://github.com/bigbio/quantms-containers} (see Experimental Procedures).
Pipeline version v2.1.0 (release codename ``Sao Pablo'', 2026-05-08) was used for the analyses reported in this paper.
Analysis scripts that generated the figures in this manuscript are available at \url{https://github.com/bigbio/quantmsdiann-paper} (this repository).

\subsection*{Competing interests}
\placeholder{VD is the original author of DIA-NN. Other authors declare no competing interests.}

\subsection*{Funding}
\placeholder{Grant numbers and funding sources for each author.}

\subsection*{Authors' contributions}
\placeholder{QXY, YS and CD developed the pipeline. AL and HW implemented the SDRF parsing and QPX export. JN and FC contributed the spatial-proteomics datasets and analysis. VD provided DIA-NN expertise and guidance on cross-version reproducibility. MB and YPR conceived and supervised the project. QXY and YPR wrote the manuscript with input from all authors. All authors read and approved the final manuscript.}

\subsection*{Acknowledgements}
\placeholder{We thank the Brunner, Derks, Leduc, Mund, Wang, and Coscia teams for depositing their data publicly. We thank the nf-core community and the OpenMS and DIA-NN developer teams for their ongoing work.}


\bibliographystyle{elsarticle-num}
\bibliography{references}

\end{document}
